\MainChapter{Conclusiones}{\topBanner}\label{cap:conclusiones}
\vspace*{0.7cm}

El desarrollo del proyecto permitió especificar, diseñar e implementar una solución de servicio de directorio orientada a centralizar la gestión de identidades y el control de acceso a los servicios considerados en el entorno del Grupo de Investigación en Redes, Información y Distribución (GRID). A partir del análisis de las necesidades identificadas y de la evaluación de las tecnologías consideradas, se seleccionó FreeIPA como la alternativa que presentó mayor correspondencia con los criterios establecidos para el proyecto.

El proceso de identificación y selección permitió establecer un conjunto de criterios relacionados con la gestión centralizada de identidades, autenticación, administración de usuarios y grupos, políticas de seguridad, compatibilidad e interoperabilidad. La aplicación del análisis DAR permitió realizar la selección de manera estructurada y fundamentada, proporcionando un mecanismo para relacionar las necesidades identificadas con las capacidades de las tecnologías evaluadas.

El diseño de la solución permitió definir una arquitectura basada en un servicio de directorio centralizado, complementado mediante un servidor réplica. La implementación de esta arquitectura permitió disponer de una copia sincronizada de la información del directorio y establecer una base para mejorar la disponibilidad del servicio. De esta manera, se redujo la dependencia de una única instancia del servicio de directorio y se establecieron las condiciones necesarias para evaluar su comportamiento ante escenarios de indisponibilidad.

La implementación de FreeIPA permitió centralizar la administración de usuarios y grupos, así como establecer políticas de seguridad para las cuentas administradas. La integración mediante LDAP permitió que los servicios contemplados en el proyecto utilizaran las identidades administradas desde un único punto, evitando mantener de manera independiente las credenciales de los usuarios en cada servicio. Asimismo, la utilización de grupos permitió establecer mecanismos de autorización asociados a los servicios, de acuerdo con la pertenencia de cada usuario.

Las pruebas realizadas permitieron verificar el funcionamiento de la solución y determinar el grado de cumplimiento de los requisitos establecidos. Los resultados evidenciaron que la solución satisface las necesidades relacionadas con la centralización de identidades, la administración de usuarios y grupos, la aplicación de políticas de seguridad, la integración con los servicios y la disponibilidad proporcionada mediante la replicación del directorio.

Durante la validación también se identificaron limitaciones relacionadas con algunas funcionalidades planteadas inicialmente. En particular, la autenticación multifactor no pudo implementarse mediante las integraciones LDAP realizadas, debido a las características de los mecanismos de autenticación utilizados por los servicios. De igual manera, la recuperación autónoma de contraseñas por parte de los usuarios no fue incorporada en la solución implementada. Estas limitaciones permitieron delimitar las capacidades efectivamente alcanzadas y evidenciaron que la incorporación de dichas funcionalidades requeriría mecanismos adicionales a la integración LDAP utilizada.

En términos generales, el proyecto permitió obtener una solución funcional de servicio de directorio que responde a las necesidades de centralización de identidades y control de acceso identificadas para el entorno evaluado. Los resultados obtenidos permiten concluir que FreeIPA constituye una alternativa viable para este escenario, siempre que las funcionalidades requeridas sean compatibles con los mecanismos de integración utilizados y que aquellas capacidades no cubiertas mediante LDAP sean abordadas mediante componentes complementarios cuando formen parte de futuras ampliaciones de la solución.

